\par\medskip\noindent
\underline{Definición 1.2. Un grupoide cofibrado $\underline A$ sobre
$\underline C_\Lambda$ se llama semihomogéneo si se cumplen las
propiedades siguientes:}

(1) \underline{Todo diagrama}
\[
\begin{tikzcd}[row sep=large,column sep=large]
a'' \arrow[r,dashed] \arrow[d,dotted] & a' \arrow[d] \\
a_1 \arrow[r] & a
\end{tikzcd}
\]
\underline{de flechas continuas en $\underline A$, donde
$a'\longrightarrow a$ es sobreyectivo, puede completarse en un diagrama
conmutativo añadiendo las flechas punteadas.}

(2) \underline{Sean}
\[
\begin{tikzcd}[column sep=large]
& R \\
R\underset{k}{\times}k[\varepsilon]
 \arrow[ur,"\pi_1"] \arrow[dr,"\pi_2"'] & \\
& k[\varepsilon]
\end{tikzcd}
\]
\underline{los morfismos de proyección, donde $k[\varepsilon]$ es el
anillo de números duales sobre $k$, y sean}
\[
\begin{tikzcd}[column sep=huge]
a' \arrow[dr,"\alpha"] & \\
& a \\
b' \arrow[ur,"\beta"'] &
\end{tikzcd}
\]
\underline{$\pi_1$-morfismos en $\underline A$. Entonces existe un
$\underline C_\Lambda$-morfismo $\gamma:a'\longrightarrow b'$ tal que
$\alpha=\beta\gamma$ si y solo si
$(\pi_2)_*(a')\simeq(\pi_2)_*(b')$.}

\underline{Observación 1.3(a).} Más adelante, en 2.6(b), daremos una
reformulación de esta noción en términos de $2$-categorías. Si
$\underline A$ es semihomogéneo, entonces $[\underline A]$ también lo
es. En particular, 1.2(2) implica que el morfismo canónico

%%%%%%%%%%  scan idx 46  |  folio 35  %%%%%%%%%%

\[
 [\underline A(R\times_k k[\varepsilon])]
 \longrightarrow
 [\underline A(R)]\times[\underline A(k[\varepsilon])]
\]
es biyectivo.

\underline{Observación 1.3(b).} Consideremos un diagrama conmutativo
\[
\begin{tikzcd}[row sep=large,column sep=large]
& a' \arrow[dl,"v_1"'] \arrow[dr,"v_2"] & \\
a_1 \arrow[dr,"u_1"'] & & a_2 \arrow[dl,"u_2"] \\
& a &
\end{tikzcd}
\]
en una categoría cofibrada $\underline A$ sobre
$\underline C_\Lambda$. Pongamos $R=P(a)$, $R'=P(a')$ y
$R_i=P(a_i)$ para $i=1,2$. Si $a''$ es la imagen directa por el morfismo
\[
 (P(v_1),P(v_2)):R'\longrightarrow R_1\times_RR_2,
\]
obtenemos un diagrama conmutativo
\[
\begin{tikzcd}[row sep=large,column sep=large]
& a' \arrow[d] \arrow[ddl,"v_1"'] \arrow[ddr,"v_2"] & \\
& a'' \arrow[dl] \arrow[dr] & \\
a_1 \arrow[dr,"u_1"'] & & a_2 \arrow[dl,"u_2"] \\
& a &
\end{tikzcd}
\]
Por tanto, 1.2(1) puede sustituirse por:

(1)$'$ Todo diagrama
\[
\begin{tikzcd}[cramped,column sep=small,row sep=small]
a_1 \arrow[dr] & & a_2 \arrow[dl] \\
& a &
\end{tikzcd}
\qquad\text{en }\underline A\text{ situado sobre}\qquad
\begin{tikzcd}[cramped,column sep=small,row sep=small]
R_1 \arrow[dr] & & R_2 \arrow[dl] \\
& R &
\end{tikzcd}
\quad\text{en }\underline C_\Lambda,
\]
donde $R_2\longrightarrow R$ es sobreyectivo, puede completarse en un
diagrama conmutativo
\[
\begin{tikzcd}[row sep=large,column sep=large]
& a' \arrow[dl] \arrow[dr] & \\
a_1 \arrow[dr] & & a_2 \arrow[dl] \\
& a &
\end{tikzcd}
\qquad\text{situado sobre}\qquad
\begin{tikzcd}[row sep=large,column sep=large]
& R_1\times_RR_2 \arrow[dl,"\pi_1"'] \arrow[dr,"\pi_2"] & \\
R_1 \arrow[dr] & & R_2 \arrow[dl] \\
& R &
\end{tikzcd}
\]
donde $\pi_i$ ($i=1,2$) son las proyecciones.

\underline{Observación 1.3(c).} Si $\underline A$ es un grupoide
cofibrado semihomogéneo sobre $\underline C_\Lambda$, entonces
$[\underline A(k[\varepsilon])]$ está provisto canónicamente de una
estructura de espacio vectorial sobre $k$, donde $k[\varepsilon]$ es el
anillo de números duales sobre $k$. En efecto, sea
$\underline{\underline V}$ la categoría de los

%%%%%%%%%%  scan idx 47  |  folio 36  %%%%%%%%%%

espacios vectoriales de dimensión finita sobre $k$. Todo funtor que
conserva los productos
$H:\underline{\underline V}\longrightarrow(\mathrm{Ens})$ se factoriza
canónicamente como
\[
\begin{tikzcd}[row sep=large,column sep=large]
\underline{\underline V} \arrow[rr,"H"]
 \arrow[dr,"\overline H"'] & & (\mathrm{Ens}) \\
& \underline{\underline V} \arrow[ur] &
\end{tikzcd}
\]
donde $\underline{\underline V}\longrightarrow(\mathrm{Ens})$ es el
funtor de olvido. Por otra parte, el funtor
\[
 G:\underline{\underline V}\longrightarrow\underline C_\Lambda,
 \qquad W\longmapsto k[W]=k\oplus W,\qquad W^2=0,
\]
lleva los productos de $\underline{\underline V}$ a productos fibrados
sobre $k$ en $\underline C_\Lambda$. Si $\underline A$ es un grupoide
cofibrado semihomogéneo sobre $\underline C_\Lambda$, entonces
\[
 [\underline A(k[V]\times_k k[W])]
 \longrightarrow
 [\underline A(k[V])]\times[\underline A(k[W])]
\]
es biyectivo por 1.3(a). Por consiguiente, el funtor
$V\longmapsto[\underline A(k[V])]$ conmuta con los productos, lo que
demuestra la afirmación.
